\newpage
\pdfbookmark[1]{\infoname}{info}

% Eesti keeles info.
\begin{info}
\begin{abstract}
Käesolevas magistritöös käsitletakse lamapuidu automaatset tuvastamist aerolaserskaneerimise andmetest ja masinnägemise meetoditega. Lamapuit on oluline metsade süsinikuringe, struktuurse mitmekesisuse ja elurikkuse indikaator, kuid selle ruumiline hindamine põhineb seni valdavalt välitöödel, mis on ajamahukad ja ruumiliselt piiratud. Töö eesmärk oli luua ja hinnata korratav metoodiline töövoog, mis ühendab LiDAR-punktipilvede eeltöötluse, madala taimkatte kõrgusmudelite loomise, märgistatud treeningandmestiku koostamise, ruumiliselt kontrollitud valideerimise ning masinnägemise mudelite võrdlemise.

Andmestiku loomisel kasutati Maa- ja Ruumiameti aerolaserskaneerimise andmeid aastatest 2017--2024. Punktipilvedest arvutati maapinnast normaliseeritud kõrgused ning loodi 0--1,3 m kõrgusvahemikku kirjeldavad 20 cm pikslisuurusega madala taimkatte kõrgusmudelid. Klassifitseerimise töövoos hõlmati 119 LAZ-faili ja 29 kaardilehte; aktiivõppega toetatud märgistamise tulemusena moodustus 580\,136 paanist koosnev andmestik. Andmelekke vähendamiseks rakendati ruumilis-ajalist train/val/test jaotust. Lõplik neljamudeline klassifitseerimisansambel saavutas sõltumatul testandmestikul F1-skoori 0,9819 ja ROC AUC väärtuse 0,9885.

Lisaks paanipõhisele klassifitseerimisele katsetati semantilist segmenteerimist Ruhnu saare eraldi märgistatud kaardilehel. Parim segmenteerimiskonfiguratsioon saavutas lukustatud testkomplektil F1-skoori 0,5750, \textit{clDice} väärtuse 0,4173 ja IoU väärtuse 0,4035. Tulemused näitavad, et LiDAR-põhine lamapuidu paanipõhine tuvastamine on Eesti oludes tehniliselt teostatav, kuid pikslitasandi kuju taastamine ja sõltumatu ökoloogiline tõlgendamine vajavad edasist arendust. Töö peamine panus on taaskasutatav andme- ja valideerimisraamistik, mida saab kasutada tulevaste ALS-andmete, metsaseire ja looduse taastamise monitooringu metoodikate arendamisel.
\end{abstract}

\keywords{lamapuit, LiDAR, masinnägemine, metsade süsinik, kaugseire}

\cercs{P176 Tehisintellekt; P170 Arvutiteadus, arvutusmeetodid, süsteemid}
\end{info}


% Inglise keeles info.
\begin{otherInfo}{english}{Automated Detection of Downed Deadwood from Airborne LiDAR Data}
\begin{abstract}
This master's thesis investigates the automated detection of downed deadwood from airborne laser scanning data using machine vision methods. Downed deadwood is an important indicator of forest carbon cycling, structural diversity, and biodiversity, but its spatial assessment still relies mainly on field surveys, which are time-consuming and spatially limited. The aim of this study was to develop and evaluate a reproducible methodological workflow combining LiDAR point cloud preprocessing, low-vegetation canopy height model generation, labelled training data creation, spatially controlled validation, and machine vision model comparison.

The dataset was built from airborne laser scanning data collected by the Estonian Land and Spatial Development Board in 2017--2024. Height-above-ground values were derived from the point clouds, and 20 cm low-vegetation height models representing the 0--1.3 m height range were generated. The classification workflow covered 119 LAZ files and 29 map sheets; active-learning-assisted annotation produced a dataset of 580\,136 image tiles. A spatial-temporal train/validation/test split was applied to reduce data leakage. The final four-model classification ensemble achieved an F1 score of 0.9819 and a ROC AUC of 0.9885 on an independent test set.

In addition to tile-level classification, semantic segmentation was tested on a separately annotated map sheet from Ruhnu Island. The best segmentation configuration achieved an F1 score of 0.5750, a \textit{clDice} score of 0.4173, and an IoU of 0.4035 on a locked test set. The results indicate that tile-level LiDAR-based detection of visible downed deadwood is technically feasible under Estonian conditions, while pixel-level shape reconstruction and independent ecological interpretation require further development. The main contribution of the thesis is a reusable data and validation framework for future airborne laser scanning datasets, forest monitoring, and nature restoration monitoring workflows.
\end{abstract}

\keywords{downed deadwood, LiDAR, machine vision, forest carbon, remote sensing}

\cercs{P176 Artificial Intelligence; P170 Computer Science, Computational Methods}
\end{otherInfo}
